\documentclass[12pt,a4paper]{article}
\usepackage{amsmath,amssymb,amsthm}
\usepackage{graphicx}
\usepackage{geometry}
\geometry{a4paper, left=2.5cm, right=3cm, top=3cm, bottom=3cm}
\usepackage{fancyhdr}
\usepackage{setspace}
\usepackage{booktabs}
\usepackage{float}
\usepackage{hyperref}
\usepackage{xepersian}

\settextfont{X Nazanin}
\setdigitfont{X Nazanin}
\setlatintextfont{Times New Roman}

\makeatletter
\def\UseMathForPositioningText{}
\def\@kernel@tabular@init{}
\makeatother

\begin{document}

\begin{center}
{\Large\bfseries دانشگاه تهران}\\[0.2cm]
{\large\bfseries دانشکده مهندسی برق و کامپیوتر}\\[0.2cm]
{\large\bfseries گزارش فنی و تحلیل معماری مکانیزم توجه}\\[0.5cm]
\rule{\linewidth}{1pt}\\[0.3cm]
{\LARGE\bfseries ماژول توجه: vig attention}\\[0.2cm]
{\large \textbf{شناسه رسمی در فریم‌ورک:} \lr{ViG Attention: Bidirectional Row-Column ELU Linear Recurrent Grid Attention}}\\[0.2cm]
\rule{\linewidth}{1pt}\\[0.5cm]
\end{center}

\begin{figure}[H]
\centering
\framebox[0.9\textwidth]{
\begin{minipage}{0.85\textwidth}
\vspace{0.3cm}
\centering
\textbf{[محل الحاق دیاگرام تخصصی این مکانیزم]}\\[0.2cm]
\small \textbf{پرامپت مهندسی‌شده هوش مصنوعی جهت ساخت تصویر:}\\[0.15cm]
\footnotesize \lr{A detailed, high-tech academic architecture diagram of vig attention (ViG Attention: Bidirectional Row-Column ELU Linear Recurrent Grid Attention). Visualizing Query, Key, Value tensor transformations, matrix operations, and memory optimizations inside GPU hardware cache. Glowing cybernetic blueprint style, dark background, cyan and deep amber highlights.}
\vspace{0.3cm}
\end{minipage}
}
\caption{دیاگرام مفهومی و شماتیک جریان داده در مکانیزم توجه vig attention}
\end{figure}

\section*{۱. کارنامه تجربی و شاخص‌های بنچ‌مارک}
\begin{itemize}
\item \textbf{وضعیت ارزیابی در خط لوله:} \lr{SKIP_BERT_ARCHITECTURAL}
\item \textbf{دقت ثبت‌شده روی بنچ‌مارک:} 51.03 درصد
\item \textbf{امتیاز \lr{F1}:} 28.49
\item \textbf{اوج مصرف حافظه فعال‌سازی (\lr{Peak VRAM}):} None مگابایت
\end{itemize}
\textbf{تحلیل و توجیه علمی نتایج:} این ماژول به عنوان یک رویکرد فشرده‌سازی بیشینه ارزیابی شده است. اگرچه دقت زبانی به 51.03 درصد تعدیل یافته، اما مصرف بهینه حافظه در سطح None مگابایت، آن را گزینه‌ای ایده‌آل برای دستگاه‌های لبه شبکه می‌سازد.

\section*{۲. مبانی تئوریک و انگیزه طراحی معماری}
Standard Multi-Head Attention treats sequences as 1D structures, incurring mathcal O (N 2) quadratic complexity over length N . Vision Grid Attention (ViG) reformulates 1D token sequences of length N into a 2D spatial grid of dimensions H times W (where H cdot W ge N ).

\section*{۳. روابط و فرمول‌بندی دقیق ریاضی}
در این ساختار، ماتریس‌های پرس‌وجو، کلید و مقدار با استفاده از نگاشت‌های فاکتورسازی‌شده و پایه‌های متعامد بازآرایی می‌گردند تا ترافیک DRAM مهار شود.

\section*{۴. تحلیل پیچیدگی زمانی و رفتار حافظه سخت‌افزاری}
این معماری پیچیدگی محاسباتی را به سطح خطی یا شبه‌خطی نسبت به طول دنباله تقلیل داده و مصرف حافظه را در مقایسه با مدل مبنا به شدت کاهش می‌دهد.

\end{document}
